% ===========================================================================
% figs/licensing.tex -- fig:licensing, sec:methods-probe.
% Two stages: (a) the sample and the tap, (b) fitting one layer's probe (the
% 5-fold split, the logistic fit itself, the readout). The outcome-asymmetry
% argument and the cell-line control are made in text, not drawn here.
% ===========================================================================
\begin{tikzpicture}[x=1cm, y=1cm, font=\small, inner sep=0pt]
  \tikzset{
    ttl/.style ={font=\footnotesize\bfseries, text=ink, anchor=base west},
    sub/.style ={font=\scriptsize\itshape, text=quietink, anchor=base west},
    lab/.style ={font=\scriptsize, text=ink, anchor=base west},
    qlab/.style={font=\scriptsize, text=quietink, anchor=base west},
    qlabe/.style={font=\scriptsize, text=quietink, anchor=base east},
    arw/.style ={-{Latex[length=3.2pt,width=2.6pt,line width=0pt]}},
  }

  % --- reused token-cell primitives (figs/hook.tex idiom) -------------------
  \newcommand{\tokbox}[4]{%
    \draw[rulegrey, line width=0.30pt, fill=panelbg, rounded corners=1pt]
      (#1,#3) rectangle (#2,{#3+0.42});
    \node[font=\scriptsize, text=ink] at ({(#1+#2)/2},{#3+0.20}) {#4};}
  \newcommand{\tokdrug}[3]{%
    \draw[accent, line width=0.40pt, fill=accent!12, rounded corners=1pt]
      (#1,#3) rectangle (#2,{#3+0.42});
    \node[font=\scriptsize, text=accent] at ({(#1+#2)/2},{#3+0.20}) {drug};}

% ###########################################################################
% # (a) the sample and the tap
% ###########################################################################
  \node[ttl] at (0.00,7.15) {(a) the sample and the tap};
  \node[sub] at (0.00,6.75)
    {$12$ drugs $\times$ $40$ cells, pooled without regard to well, $480$ prompts};

  % --- the stack: seven residual-stream taps, one per probed depth ---------
  \def\lx{0.55} \def\lw{1.90} \def\ly{1.35} \def\lh{0.52} \def\gap{0.14}
  \foreach \i/\name/\hi in {0/2/0,1/4/0,2/6/0,3/8/0,4/9/1,5/12/0,6/16/0}{
    \pgfmathsetmacro{\yb}{\ly + \i*(\lh+\gap)}
    \ifnum\hi=1
      \draw[accent, line width=0.55pt, fill=accent!18, rounded corners=1pt]
        (\lx,\yb) rectangle ({\lx+\lw},{\yb+\lh});
      \node[font=\scriptsize, text=accent] at ({\lx+\lw/2},{\yb+\lh/2-0.05})
        {layer \name};
    \else
      \draw[rulegrey, line width=0.30pt, fill=panelbg, rounded corners=1pt]
        (\lx,\yb) rectangle ({\lx+\lw},{\yb+\lh});
      \node[font=\scriptsize, text=quietink] at ({\lx+\lw/2},{\yb+\lh/2-0.05})
        {layer \name};
    \fi
  }
  \node[qlabe, rotate=90, anchor=south]
    at ({\lx-0.18},{\ly+3.5*(\lh+\gap)}) {depth (index, not to scale)};

  % --- the prompt strip, five fields, drug tinted (figs/hook.tex order) ----
  \def\py{-1.00}
  \tokbox{0.00}{1.55}{\py}{cell line}
  \tokdrug{1.55}{2.75}{\py}
  \tokbox{2.75}{3.75}{\py}{dose}
  \tokbox{3.75}{5.35}{\py}{mechanism}
  \tokbox{5.35}{7.85}{\py}{control cell sentence}
  \node[qlab] at (0.00,{\py-0.42}) {the prompt, five fields (\cref{sec:methods-data})};
  \draw[rulegrey, line width=0.35pt, arw] ({\lx+\lw/2},{\py+0.44}) -- ({\lx+\lw/2},{\ly-0.02});

  % --- the tap: layer 9's residual stream, pulled out ------------------------
  \pgfmathsetmacro{\taprow}{\ly + 4*(\lh+\gap) + \lh/2}
  \draw[accent, line width=0.7pt, arw] ({\lx+\lw},\taprow) -- (10.30,\taprow);
  \node[lab, text=accent, anchor=south west]
    at ({\lx+\lw+0.10},{\taprow+0.10}) {activation vector};
  \node[qlab, anchor=north west]
    at ({\lx+\lw+0.10},{\taprow-0.10}) {2048-wide, one position, one of seven independent reads};
  \node[qlab, anchor=west] at (10.45,\taprow) {carried down to (b) $\downarrow$};
  \draw[accent, line width=0.5pt, arw] (10.30,{\taprow-0.20}) -- (10.30,-1.05);

% ###########################################################################
% # (b) fitting one layer's probe
% ###########################################################################
  \def\by{-2.55}
  \draw[rulegrey, line width=0.30pt] (0.00,{\by+0.75}) -- (17.00,{\by+0.75});
  \node[ttl] at (0.00,{\by+0.30}) {(b) fitting one layer's probe};
  \node[sub] at (0.00,{\by-0.10})
    {a separate logistic regression at each depth, none sharing weights with another};

  % --- the split: 480 prompts, five folds, one held out ---------------------
  \def\sx{0.00} \def\sy{\by-0.85} \def\sw{5.60} \def\sh{0.48}
  \node[lab] at (\sx,{\sy+0.28}) {$480$ prompts, $5$ folds};
  \foreach \k/\hi in {0/0,1/0,2/0,3/0,4/1}{
    \pgfmathsetmacro{\xa}{\sx + \k*\sw/5}
    \pgfmathsetmacro{\xb}{\sx + (\k+1)*\sw/5}
    \ifnum\hi=1
      \draw[accent, line width=0.55pt, fill=accent!15, rounded corners=1pt]
        (\xa,\sy) rectangle (\xb,{\sy+\sh});
    \else
      \draw[rulegrey, line width=0.30pt, fill=panelbg, rounded corners=1pt]
        (\xa,\sy) rectangle (\xb,{\sy+\sh});
    \fi
  }
  \node[qlab, anchor=north west] at (\sx,{\sy-0.14})
    {stratified on the drug label};
  \node[qlab, text=accent, anchor=north west] at (\sx,{\sy-0.50})
    {\parbox{5.7cm}{fit on the four grey folds, scored on
     the held-out accent fold --- shown once; repeated per
     fold, the five scores averaged}};

  % --- the fit: standardise, then a 12-way logistic layer --------------------
  \def\mx{6.55} \def\mh{0.62}
  \pgfmathsetmacro{\my}{\sy + \sh/2 - \mh/2}
  \draw[ink, line width=0.5pt, arw] (\sx+\sw+0.10,{\sy+\sh/2}) -- (\mx,{\my+\mh/2});
  \draw[ink, line width=0.45pt, fill=panelbg, rounded corners=1pt]
    (\mx,\my) rectangle ({\mx+1.75},{\my+\mh});
  \node[font=\scriptsize, text=ink] at ({\mx+0.875},{\my+\mh/2})
    {\parbox{1.55cm}{\centering standardise}};
  \draw[ink, line width=0.5pt, arw] ({\mx+1.75},{\my+\mh/2}) -- ({\mx+2.15},{\my+\mh/2});
  \draw[ink, line width=0.55pt, fill=accent!10, rounded corners=1pt]
    ({\mx+2.15},\my) rectangle ({\mx+4.55},{\my+\mh});
  \node[font=\scriptsize, text=ink]
    at ({\mx+3.35},{\my+\mh/2})
    {\parbox{2.15cm}{\centering logistic\\(12-way)}};
  \node[qlab, anchor=north west] at (\mx,{\my-0.20})
    {\parbox{4.4cm}{one weight vector per drug; the class with
     the largest score wins ($C=1.0$, untuned)}};

  % --- the readout: twelve candidate drugs, chance line, printed accuracy --
  \def\ox{11.60} \def\ow{5.30} \def\oy{\my+\mh/2}
  \draw[ink, line width=0.5pt, arw] ({\mx+4.55},{\my+\mh/2}) -- (\ox,\oy);
  \foreach \k in {0,...,11}{
    \pgfmathsetmacro{\xx}{\ox + \k*\ow/11}
    \ifnum\k=3
      \fill[accent] (\xx,\oy) circle (2.0pt);
    \else
      \fill[rulegrey] (\xx,\oy) circle (1.1pt);
    \fi
  }
  \draw[ink, line width=0.35pt, densely dotted]
    (\ox,{\oy-0.32}) -- ({\ox+\ow},{\oy-0.32});
  \node[qlabe] at ({\ox+\ow},{\oy-0.32-0.24}) {chance $=1/12\approx 0.083$};
  \pgfmathsetmacro{\accx}{\ox + 3*\ow/11}
  \node[font=\scriptsize, text=accent] at (\accx,{\oy+0.30}) {$0.821$};
  \node[qlab, anchor=south west] at (\ox,{\oy+0.55})
    {\parbox{5.3cm}{twelve candidates, five folds averaged; accent
     is the true drug}};

% ###########################################################################
% # (c) accuracy holds up; the geometry does not
% ###########################################################################
% Source: docs/endcell/dimensionality_probe_analysis.md, the retrained
% [END_CELL] model's table -- every layer 0-16, not only the seven tapped in
% (a)/(b). Two quantities, two scales, one shared x (layer index), following
% the same convention as fig:sweepbins: the divider states it rather than a
% dual axis implying a false common scale. Accent marks the seven layers (a)
% and (b) already tap; quietink marks the other ten, present here because this
% instrument reads every layer and (a)/(b)'s does not.
  \def\cy{\by-3.85}
  \draw[rulegrey, line width=0.30pt] (0.00,{\cy+0.75}) -- (17.00,{\cy+0.75});
  \node[ttl] at (0.00,{\cy+0.30}) {(c) accuracy holds up; the geometry does not};
  \node[sub] at (0.00,{\cy-0.10})
    {every layer, $0$ to $16$ --- a denser sweep than the seven tapped above};

  \def\cx{0.90} \def\cstep{0.90}
  \newcommand{\clx}[1]{\pgfmathresult}

  % --- upper strip: probe accuracy, one dot per layer -----------------------
  \def\ay{\cy-3.50}
  \node[lab] at (0.00,{\ay+2.85}) {probe accuracy};
  \draw[ink, line width=0.35pt, densely dotted]
    ({\cx},{\ay+0.234}) -- ({\cx+16*\cstep},{\ay+0.234});
  \node[qlabe] at ({\cx+16*\cstep},{\ay+0.234-0.24}) {chance $\approx 0.08$};
  \draw[quietink, line width=0.35pt]
    plot coordinates {
     (0.90,{\ay+0.234}) (1.80,{\ay+0.912}) (2.70,{\ay+1.152}) (3.60,{\ay+1.499})
     (4.50,{\ay+1.652}) (5.40,{\ay+1.841}) (6.30,{\ay+1.954}) (7.20,{\ay+2.053})
     (8.10,{\ay+2.200}) (9.00,{\ay+2.318}) (9.90,{\ay+2.211}) (10.80,{\ay+2.112})
     (11.70,{\ay+2.129}) (12.60,{\ay+2.106}) (13.50,{\ay+2.171}) (14.40,{\ay+2.160})
     (15.30,{\ay+2.140})};
  \foreach \L/\x/\v/\tap in {0/0.90/0.234/0,1/1.80/0.912/0,2/2.70/1.152/1,
     3/3.60/1.499/0,4/4.50/1.652/1,5/5.40/1.841/0,6/6.30/1.954/1,7/7.20/2.053/0,
     8/8.10/2.200/1,9/9.00/2.318/1,10/9.90/2.211/0,11/10.80/2.112/0,
     12/11.70/2.129/1,13/12.60/2.106/0,14/13.50/2.171/0,15/14.40/2.160/0,
     16/15.30/2.140/1}{
    \ifnum\tap=1
      \fill[accent] (\x,{\ay+\v}) circle (1.6pt);
    \else
      \fill[quietink] (\x,{\ay+\v}) circle (0.9pt);
    \fi
  }
  \node[font=\scriptsize, text=accent] at (9.00,{\ay+2.318+0.30}) {$0.821$};
  \node[font=\scriptsize, text=accent] at (15.30,{\ay+2.140+0.30}) {$0.758$};

  % --- divider: two quantities, two scales, one x ---------------------------
  \draw[rulegrey, line width=0.30pt] (0.00,{\ay-0.55}) -- (17.00,{\ay-0.55});

  % --- lower strip: between/within variance ratio, one dot per layer --------
  \def\ry{\cy-6.85}
  \node[lab] at (0.00,{\ry+2.40})
    {between/within variance ratio, same activations};
  \draw[ink, line width=0.35pt] ({\cx},{\ry+0.00}) -- ({\cx+16*\cstep},{\ry+0.00});
  \node[qlabe] at ({\cx+16*\cstep},{\ry-0.24}) {$0$: no separation};
  \draw[quietink, line width=0.35pt]
    plot coordinates {
     (1.80,{\ry+1.067}) (2.70,{\ry+1.117}) (3.60,{\ry+1.433}) (4.50,{\ry+1.717})
     (5.40,{\ry+1.717}) (6.30,{\ry+1.800}) (7.20,{\ry+1.783}) (8.10,{\ry+1.817})
     (9.00,{\ry+1.867}) (9.90,{\ry+1.400}) (10.80,{\ry+1.367}) (11.70,{\ry+0.700})
     (12.60,{\ry+0.583}) (13.50,{\ry+0.583}) (14.40,{\ry+0.567}) (15.30,{\ry+0.567})};
  \foreach \L/\x/\v/\tap in {1/1.80/1.067/0,2/2.70/1.117/1,3/3.60/1.433/0,
     4/4.50/1.717/1,5/5.40/1.717/0,6/6.30/1.800/1,7/7.20/1.783/0,8/8.10/1.817/1,
     9/9.00/1.867/1,10/9.90/1.400/0,11/10.80/1.367/0,12/11.70/0.700/1,
     13/12.60/0.583/0,14/13.50/0.583/0,15/14.40/0.567/0,16/15.30/0.567/1}{
    \ifnum\tap=1
      \fill[accent] (\x,{\ry+\v}) circle (1.6pt);
    \else
      \fill[quietink] (\x,{\ry+\v}) circle (0.9pt);
    \fi
  }
  \node[font=\scriptsize, text=accent] at (9.00,{\ry+1.867+0.30}) {$0.112$};
  \node[font=\scriptsize, text=accent] at (15.30,{\ry+0.567-0.30}) {$0.034$};

  % --- shared x-axis: layer index --------------------------------------------
  \foreach \L/\x in {0/0.90,4/4.50,8/8.10,9/9.00,12/11.70,16/15.30}{
    \node[font=\scriptsize, text=quietink, anchor=north]
      at (\x,{\ry-0.55}) {\L};}
  \node[qlab] at (7.20,{\ry-0.95}) {layer};

\end{tikzpicture}
